\begin{table}[H]
\centering
\caption{Dollaren som safe haven: ugeafkast i stressuger før og efter april 2025}
\label{tab:campbell_safehaven}
\begin{threeparttable}
\begin{tabular}{l ccc}
\toprule
 & Før april 2025 & Efter april 2025 & Forskel \\
\midrule
\multicolumn{4}{l}{\textit{Panel A: stressuge = nederste 10 pct. af aktieafkastene}} \\
Hældning, normale uger (hedge) & $-0.17^{***}$ & $-0.18^{**}$ & $-0.01$ \\
 & $(0.03)$ & $(0.09)$ & $(0.09)$ \\
Hældning, stressuger (safe haven) & $-0.03$ & $0.08$ & $0.11$ \\
 & $(0.07)$ & $(0.21)$ & $(0.22)$ \\
Gns. dollarafkast i stressuger (pct.) & $0.10$ & $0.17$ & $0.07$ \\
 & $(0.18)$ & $(0.73)$ & $(0.76)$ \\
Antal stressuger & 57 & 5 & \\[4pt]
\multicolumn{4}{l}{\textit{Panel B: stressuge = nederste 5 pct. af aktieafkastene}} \\
Hældning, normale uger (hedge) & $-0.13$ & $-0.16$ & $-0.03$ \\
 & $(nan)$ & $(nan)$ & $(nan)$ \\
Hældning, stressuger (safe haven) & $-0.07$ & $0.13$ & $0.20$ \\
 & $(nan)$ & $(nan)$ & $(nan)$ \\
Gns. dollarafkast i stressuger (pct.) & $0.38$ & $-1.13^{*}$ & $-1.51^{**}$ \\
 & $(0.24)$ & $(0.67)$ & $(0.71)$ \\
Antal stressuger & 30 & 1 & \\[4pt]
\bottomrule
\end{tabular}
\begin{tablenotes}[flushleft]\footnotesize
\item \textit{Note:} Ugentlige merafkast. Hældningerne er fra én regression af dollarens merafkast over for euroen på de amerikanske aktiers merafkast (S\&P 500), opdelt på normale uger og stressuger før og efter april 2025 (Baur og Lucey 2010). En negativ hældning betyder, at dollaren styrkes, når aktierne falder. En stressuge er en uge, hvor aktiernes merafkast ligger i den nederste fraktil for hele stikprøven. HC3-standardfejl i parentes. $^{*}$, $^{**}$ og $^{***}$ angiver signifikans på 10, 5 og 1 pct.-niveau.
\end{tablenotes}
\end{threeparttable}
\end{table}